\chapter[مقدمه]{مقدمه}

\section{بیان مسئله}

حمل‌ونقل بار جاده‌ای بخش عمده‌ای از جابه‌جایی کالا در ایران را تشکیل می‌دهد. در عمل، اتصال میان صاحب‌بار یا متصدی حمل با رانندهٔ مناسب هنوز اغلب از مسیرهای غیرساخت‌یافته انجام می‌شود: تماس تلفنی، گروه‌های پیام‌رسان، یا واسطه‌های سنتی. این روش‌ها چند مشکل مشخص ایجاد می‌کنند:

\begin{itemize}
  \item \textbf{عدم تطابق ماشین و بار:} بار یخچالی نباید به کفی معمولی پیشنهاد شود؛ نبود فیلتر نوع بارگیر باعث اتلاف وقت دو طرف است.
  \item \textbf{شفاف نبودن وضعیت:} راننده و متصدی تصویر واحدی از این‌که بار «باز»، «تخصیص‌یافته» یا «تحویل‌شده» است ندارند.
  \item \textbf{ضعف موقعیت مکانی:} پیشنهاد بار نزدیک به محل فعلی راننده، بدون موقعیت واقعی \textenglish{GPS}، عملاً تصادفی است.
  \item \textbf{مسیریابی غیردقیق:} خط مستقیم روی نقشه جایگزین مسیر خیابان و جاده می‌شود و برای راننده کاربردی نیست.
  \item \textbf{تأیید یک‌طرفهٔ اتمام کار:} اگر فقط یک طرف مأموریت را بسته اعلام کند، اختلاف حساب و وضعیت بار باقی می‌ماند.
\end{itemize}

سامانهٔ \appname\ برای پاسخ به همین خلأ طراحی شده است: یک کاربرد دوطرفه با قرارداد مشخص داده، نقش‌های جدا، چرخهٔ وضعیت بار، و مسیریابی واقعی.

\section{هدف پروژه}

هدف کلی، طراحی و پیاده‌سازی یک سامانهٔ سرتاسری (کلاینت + سرویس‌دهنده + پایگاه داده) است که:

\begin{enumerate}
  \item ثبت، ویرایش، حذف و مشاهدهٔ بار را برای متصدی ممکن کند؛
  \item فهرست بارهای باز را با فیلتر استان مبدأ/مقصد و بازهٔ قیمت در اختیار راننده بگذارد؛
  \item پذیرش بار را به راننده‌ای محدود کند که نوع ماشین او با بار سازگار باشد؛
  \item مأموریت‌های تخصیص‌یافته را جدا از بارهای باز نگه دارد؛
  \item مسیر خیابانی از موقعیت راننده تا مبدأ بار و سپس تا مقصد را نمایش دهد؛
  \item اتمام کار را منوط به تأیید راننده \emph{و} متصدی کند؛
  \item روی موبایل و به صورت وب پیش‌رونده قابل استفاده باشد.
\end{enumerate}

\section{تفاوت با سامانه‌های مشابه}

در بازار ایران سامانه‌هایی مانند \textbf{ترابرنت} و \textbf{دیان} عمدتاً برای معرفی بار و اتصال راننده به صاحب‌بار یا متصدی به کار می‌روند. تمرکز آن‌ها بر انتشار آگهی بار و پیدا کردن طرف مقابل است؛ زنجیرهٔ عملیاتی پس از پیدا شدن بار (حرکت روی جاده، تشخیص بار نزدیک به محل فعلی راننده، پذیرش رسمی، بستن دوطرفهٔ کار، و نگهداری سابقه) یا غایب است یا بیرون از خودِ برنامه باقی می‌ماند.

سامانهٔ \appname\ همین چهار حلقه را داخل کاربرد پیاده کرده است. جدول~\ref{tab:similar-apps} این تمایز را خلاصه می‌کند.

\begin{table}[ht]
\centering
\caption{چهار قابلیت تمایزدهنده نسبت به ترابرنت و دیان}
\label{tab:similar-apps}
\begin{tabularx}{\textwidth}{@{}
  >{\raggedleft\arraybackslash}p{0.22\textwidth}
  Y
  Y
@{}}
\toprule
\thcell{قابلیت} & \thcell{ترابرنت و دیان} & \thcell{سامانهٔ لجستیک} \\
\midrule
مسیریابی
  & معرفی بار و مسیر کلی؛ هدایت خیابانی زنده داخل برنامه محور کار نیست
  & مسیر واقعی با نقشهٔ نشان از موقعیت راننده تا مبدأ بار و سپس تا مقصد، با ترافیک و ناوبری \\
\midrule
بارهای نزدیک با \textenglish{GPS}
  & فهرست بار بدون فاصلهٔ لحظه‌ای تا محل راننده
  & با روشن بودن موقعیت، فاصلهٔ هر بار تا راننده روی کارت می‌آید و بار داخل شعاع نزدیک برچسب «نزدیک» می‌گیرد \\
\midrule
پذیرش و تأیید اتمام
  & اتصال غالباً با تماس یا توافق بیرون از چرخهٔ وضعیت در برنامه
  & راننده بار را در سامانه می‌پذیرد؛ اتمام کار فقط با تأیید راننده \emph{و} متصدی قطعی می‌شود \\
\midrule
سوابق و نتایج بار
  & تأکید بر بار جاری و آگهی جدید
  & راننده مأموریت‌های فعال و تکمیل‌شده را جدا می‌بیند؛ متصدی وضعیت و نتیجهٔ بارهای ثبت‌شده را در فهرست خود پیگیری می‌کند \\
\bottomrule
\end{tabularx}
\end{table}

این چهار مورد با نیاز مطرح‌شده در بیان مسئله یکی است و دلیل ارجحیت عملیاتی سامانه نسبت به ابزارهای مشابه را مشخص می‌کند:

\begin{enumerate}
  \item \textbf{مسیریابی:} پس از پذیرش، راننده به صفحهٔ مسیر می‌رود. نقشه مسیر خیابانی را تا محل بارگیری و سپس تخلیه نشان می‌دهد؛ خط مستقیم یا آدرس متنی جایگزین ناوبری نیست.
  \item \textbf{بارهای نزدیک:} با فعال‌سازی \textenglish{GPS}، فاصله تا مبدأ هر بار روی کارت نمایش داده می‌شود و بارهای داخل شعاع نزدیک در بالای فهرست جدا می‌شوند. راننده بار دور را با بار دمِ دست اشتباه نمی‌گیرد.
  \item \textbf{پذیرش و اتمام کار:} پذیرش یک عمل ثبت‌شده در سرور است و وضعیت بار را به «تخصیص‌یافته» می‌برد. بستن مأموریت یک‌طرفه ممکن نیست؛ تا هر دو طرف تکمیل را نزنند، بار «تحویل‌شده» نمی‌شود.
  \item \textbf{سوابق و نتایج:} مأموریت‌های راننده با فیلتر فعال، تکمیل‌شده و همه قابل مرور است. متصدی همان چرخه را در فهرست بارهای خود می‌بیند: در انتظار، تخصیص‌یافته، تحویل‌شده یا لغو. نتیجهٔ کار داخل سامانه می‌ماند، نه در حافظهٔ تماس‌ها.
\end{enumerate}

\section{تقسیم کار دانشجویان}

پروژه به صورت تیمی و با مرز مشخص لایه‌ها انجام شده است. جدول~\ref{tab:work-split} این تقسیم را نشان می‌دهد.

\begin{table}[ht]
\centering
\caption{تقسیم مسئولیت در پروژه}
\label{tab:work-split}
\vspace{2pt}
\begin{tikzpicture}[
  card/.style={
    rectangle split,
    rectangle split parts=2,
    rectangle split part fill={tablehead,soft},
    rectangle split draw splits=false,
    draw=brand,
    line width=0.7pt,
    rounded corners=5pt,
    inner xsep=11pt,
    inner ysep=8pt,
    minimum width=6.15cm,
  }
]
\node[card] (fe) at (3.52,0) {%
  \fanode{5.7cm}{%
    \textbf{حانیه نباتی}\\[1pt]
    {\footnotesize\textcolor{muted}{\textenglish{\ttfamily 40121160028}}}
  }
  \nodepart{two}
  \fanode{5.7cm}{%
    {\small\bfseries\textcolor{ink}{فرانت‌اند}}\\[5pt]
    رابط کاربری موبایل و وب\\
    معماری کلاینت فلاتر\\
    نقشه و وب پیش‌رونده\\
    اتصال به وب‌سرویس و آزمون
  }
};
\node[card] (be) at (-3.52,0) {%
  \fanode{5.7cm}{%
    \textbf{مهیار عامری}\\[1pt]
    {\footnotesize\textcolor{muted}{\textenglish{\ttfamily 40121160018}}}
  }
  \nodepart{two}
  \fanode{5.7cm}{%
    {\small\bfseries\textcolor{ink}{بک‌اند و پایگاه داده}}\\[5pt]
    سرویس‌دهندهٔ جنگو\\
    وب‌سرویس و لایهٔ رست\\
    مدل دامنه و پایگاه داده\\
    احراز هویت با ژتون دسترسی
  }
};
\end{tikzpicture}
\end{table}

استاد راهنمای پروژه جناب آقای \textbf{دکتر سید مرتضی بابامیر}، استاد مهندسی نرم‌افزار دانشگاه کاشان هستند.

\section{قلمرو و محدودیت‌ها}

قلمرو سامانه، حمل بار جاده‌ای داخلی با دو نقش راننده و متصدی است. پرداخت برخط، بورس بار سراسری، و ناوگان چندراننده‌ای در نسخهٔ جاری پیاده‌سازی نشده است. اعلان فشاری واقعی (\textenglish{FCM}) در کلاینت به صورت زیرساخت آماده مانده و هنوز به نقطهٔ پایانی سرور وصل نیست. پایگاه پیش‌فرض جنگو (\textenglish{SQLite3}) برای توسعه و استقرار آموزشی انتخاب شده و برای بار بسیار زیاد تولید می‌توان آن را با \textenglish{PostgreSQL} جایگزین کرد، بدون تغییر منطق \textenglish{ORM}.

\section{ساختار گزارش}

همین فصل پس از بیان مسئله و هدف، تفاوت سامانه با برنامه‌های مشابه (ترابرنت و دیان) را می‌آورد. فصل دوم نیازمندی‌ها و مدل کسب‌وکار را بیان می‌کند. فصل سوم معماری کلی را نشان می‌دهد. فصل چهارم پیاده‌سازی فرانت‌اند را شرح می‌دهد. فصل پنجم بک‌اند جنگو و \textenglish{DRF} را توضیح می‌دهد. فصل ششم پایگاه داده و نگاشت شیء-رابطه‌ای است. فصل هفتم قرارداد \textenglish{API} و امنیت است. فصل هشتم استقرار و آزمون را می‌آورد و فصل نهم جمع‌بندی است.
